\subsection{开等价关系与闭等价关系}

定义 2. 拓扑空间 X 上的等价关系 R 称为开（相应地，闭）的，如果 X 到 X/R 上的典范映射是开（相应地，闭）映射。

这等价于说，X 的每个开（相应地，闭）子集关于 R 的饱和化在 X 中是开（相应地，闭）集（§ 3，第 4 目）。

例。1) 设 X 是拓扑空间，\(\Gamma\) 是 X 到自身的一些同胚所成的一个群，并设 R 是等价关系

\[
" \text { there   exists } \sigma \in \Gamma \text { such   that } y = \sigma (x)"
\]

（x 与 y 之间）（于是 R 是以 \(\Gamma\) 在 X 中的轨道为类的那个等价关系。关系 R 是开的，因为 X 的子集 A 关于 R 的饱和化是开集，诸 \(\sigma(\mathbf{A})\) 也是开集，从而它们的并也是开集。

* 在实直线 R 上，等价关系 \(x \equiv y \pmod{1}\) 是开的，因为它如上所述由平移群

\[
x \rightarrow x + n \quad (n \in \mathbf {Z})
\]

导出（见第三章，§ 2，第 4 目）。*

\begin{enumerate}
\def\labelenumi{\arabic{enumi})}
\setcounter{enumi}{1}
\tightlist
\item
  设 X 是子空间族 \((\mathbf{X}_{t})\) 的和，并设 X/R 是把诸 \(X_{t}\) 沿开子集 \(A_{tx}\) 借助双射 \(h_{xt}\) 粘合起来所得的空间（§ 2，第 5 目）；并设 \(h_{xt}\) 是 \(A_{tx}\) 到 \(A_{xt}\) 上的一个同胚（对每一对指标 \((t, x)\)）。
\end{enumerate}

则关系 R 是开的。因为若 U 在 X 中是开集，则 U 的饱和化是诸 \(h_{xt}(\mathbf{U} \cap \mathbf{A}_{\mathbf{tx}})\) 之并；由于 \(U \cap A_{tx}\) 在 \(A_{tx}\) 中是开集，故 \(h_{xt}(\mathbf{U} \cap \mathbf{A}_{\mathbf{tx}})\) 在 \(A_{xt}\) 中从而在 X 中是开集。

\begin{enumerate}
\def\labelenumi{\arabic{enumi})}
\setcounter{enumi}{2}
\tightlist
\item
  沿用例 2 的记号，现在设诸 \(\mathbf{A}_{\mathrm{tx}}\) 是闭集且诸 \(h_{xt}\) 是同胚；再设对每个指标 \(t\)，只有有限多个指标 \(x\) 使 \(\mathbf{A}_{\mathrm{tx}} \neq \emptyset\)（亦即每个 \(\mathbf{X}_t\) 只与有限多个 \(\mathbf{X}_x\) 相“粘”）。则关系 R 是闭的。因为若 F 是 X 的任一闭子集，则 F 的饱和化是诸集合 \(h_{xt}(\mathbf{F} \cap \mathbf{A}_{\mathrm{tx}}) \subset \mathbf{A}_{xt}\) 之并；所作的假定蕴含这一族是局部有限的，且 \(h_{xt}(\mathbf{F} \cap \mathbf{A}_{\mathrm{tx}})\) 在 \(\mathbf{A}_{xt}\) 中从而在 X 中是闭集。于是结论由 § 2 第 5 目命题 4 得出。
\end{enumerate}

命题 3. 设 X、Y 是两个拓扑空间，设 \(f: \mathbf{X} \to \mathbf{Y}\) 是连续映射，设 R 是 X 上的等价关系 \(f(x) = f(y)\)，并设 \(\mathbf{X} \xrightarrow{p} \mathbf{X}/\mathbf{R} \xrightarrow{h} f(\mathbf{X}) \xrightarrow{i} \mathbf{Y}\) 是 \(f\) 的典范分解。则下列三个命题等价：

\begin{enumerate}
\def\labelenumi{\alph{enumi})}
\item
  \(f\) 是开映射。
\item
  三个映射 \(p, h, i\) 都是开映射。
\item
  等价关系 R 是开的，h 是同胚，且 \(f(\mathbf{X})\) 是 Y 的开子集。
\end{enumerate}

还有，若把前面命题中的“开”处处换为“闭”，该命题仍成立。

由于单射 \(i\) 连续，由第 1 目命题 1c) 可知，若 \(f\) 是开映射，则 \(h \circ p\) 亦然；由于 \(p\) 满且连续，命题 1b) 表明 \(h\) 是开映射；\(h\) 无论如何是连续双射；于是 \(h\) 是同胚，从而由第 1 目命题 1a)，\(p = \overline{h}^1 \circ (h \circ p)\) 是开映射。又有{[}第 1 目，命题 1b){]}，\(i \circ h\) 是开映射；故{[}第 1 目，命题 1a){]}，\(i = (i \circ h) \circ \overline{h}^1\) 亦然。这就证明了 a) 蕴含 b)。反之，第 1 目命题 1a) 表明 b) 蕴含 a)。最后，b) 与 c) 的等价性由定义立即得出。

闭映射情形的证明，作适当改动后是类似的。

命题 4. 设 R 是拓扑空间 X 上的开（相应地，闭）等价关系，f 是典范映射 X → X/R。设 A 是 X 的一个子集，且下列两个条件之一成立：a) A 在 X 中是开（相应地，闭）集。

\begin{enumerate}
\def\labelenumi{\alph{enumi})}
\setcounter{enumi}{1}
\tightlist
\item
  A 关于 R 是饱和的。
\end{enumerate}

则在 A 上诱导的关系 \(R_{A}\) 是开（相应地，闭）的，且 \(A/R_{A}\) 到 \(f(A)\) 上的典范映射是同胚。

考虑 § 3 第 6 目的交换图 (1)，它给出 \(f \circ j\) 的典范分解。在条件 a) 下，\(j\) 是开（相应地，闭）映射，而按假设 \(f\) 亦然；故 \(f \circ j\) 是开（相应地，闭）映射{[}第 1 目，命题 1 a){]}，结论由命题 3 得出。在条件 b) 下，我们有

\[
\mathbf {A} = \overline {{{f}}} ^ {1} (f (\mathbf {A})),
\]

且 \(h \circ \varphi\) 是 A 到 \(f(A)\) 中的、在 A 上与 \(f\) 一致的映射；凭借第 1 目命题 2 a)，\(h \circ \varphi\) 是开（相应地，闭）映射，而把命题 3 应用于 \(h \circ \varphi\)，又得出结论。

